\chapter{Limit dan Kekontinuan}\label{ch:b1:continuity}

Teorema nilai antara dan teorema nilai ekstrem dahulu
dipakai di tingkat Sekolah Menengah atas kepercayaan visual. Dengan barisan
(\cref{ch:b1:seq}) dan topologi $\R$ (\cref{ch:b1:topology})
di tangan, bab ini membuktikan keduanya --- lalu melengkapi teorinya dengan
teorema bijeksi monoton (yang mengesahkan $\arcsin$,
$\operatorname{arcosh}$ dan kerabatnya dari \cref{ch:b1:functions})
serta teorema Heine tentang \omterm{def:b1:continuity:uniform}{kekontinuan seragam}.

Di sepanjang bab ini, $I$ sebuah \omterm{prop:b1:reals:intervals}{selang} dan $f \colon I \to \R$; adapun ``$x_0 \in
\overline I$'' mengizinkan limit di titik ujungnya.

\section{Limit fungsi}

\begin{definition}[Limit di sebuah titik]\label{def:b1:continuity:limit}
Misalkan $x_0 \in \overline{I}$ dan $\ell \in \R$. Maka $f(x) \to \ell$ ketika
$x \to x_0$ bila\index{limit!fungsi}
\[
\forall \varepsilon > 0,\ \exists \delta > 0,\ \forall x \in I,
\qquad \abs{x - x_0} \leq \delta \implies \abs{f(x) - \ell} \leq
\varepsilon .
\]
Adapun limit di $\pm\infty$ dan limit yang tak hingga didefinisikan dengan pola yang
sama ($\abs{x - x_0} \leq \delta$ menjadi $x \geq M$; sedangkan $\abs{f(x) -
\ell} \leq \varepsilon$ menjadi $f(x) \geq M'$). Adapun limit sepihak
membatasi $x$ pada $x > x_0$ (ditulis $x \to x_0^+$) atau $x < x_0$. Dan
limitnya tunggal bila ada (dengan bukti yang sama seperti
\cref{prop:b1:seq:first}).
\end{definition}

\begin{example}[Sebuah limit lewat apitan]\label{ex:b1:continuity:floorlimit}
Hitunglah $\lim_{x \to 0} x\,\bigl\lfloor \frac1x \bigr\rfloor$.
Apitan \omterm{thm:b1:reals:floor}{lantainya} $\frac1x - 1 < \lfloor \frac1x \rfloor \leq
\frac1x$ memberikan, setelah dikalikan dengan $x$ (awas tandanya!):
\[
1 - x < x\Bigl\lfloor \frac1x \Bigr\rfloor \leq 1 \quad (x > 0),
\qquad
1 \leq x\Bigl\lfloor \frac1x \Bigr\rfloor < 1 - x \quad (x < 0),
\]
dan kedua apitan sepihaknya menutup pada $1$: jadi limitnya $1$. Perhatikan
apa yang terjadi: bahwa $\lfloor \frac1x\rfloor$ sendirian melompat liar di dekat
$0$, tetapi faktor $x$ menjinakkan setiap lompatannya (karena $x$ dikali lompatan satuan itu
kecil), dan hanya apitannya yang bertahan. Inti gagasan penutupnya: bahwa
limit hasil kali sebuah faktor yang kecil dengan faktor yang berayun terbatas
adalah soal apitan, dan tak pernah soal teorema operasi ---
sebab teorema operasi menuntut \emph{kedua} faktornya konvergen.
\end{example}

\begin{theorem}[Pencirian lewat barisan]\label{thm:b1:continuity:seqchar}
Di sini $f(x) \to \ell$ ketika $x \to x_0$ jika dan hanya jika: untuk \emph{setiap}
barisan $(u_n)$ berisi titik $I$ dengan $u_n \to x_0$, berlaku
$f(u_n) \to \ell$.
\end{theorem}

\begin{proof}
($\Rightarrow$) Misalkan $u_n \to x_0$ dan $\varepsilon > 0$. Ambillah
$\delta$ dari definisinya, lalu $N$ dengan $\abs{u_n - x_0} \leq
\delta$ untuk $n \geq N$: maka di luar $N$, $\abs{f(u_n) - \ell} \leq
\varepsilon$.

($\Leftarrow$) Lewat kontraposisinya. Jika $f \not\to \ell$: maka suatu
$\varepsilon_0 > 0$ mengalahkan setiap $\delta$; lalu memilih $\delta =
\frac{1}{n+1}$ menghasilkan $u_n \in I$ dengan $\abs{u_n - x_0} \leq
\frac{1}{n+1}$ dan $\abs{f(u_n) - \ell} > \varepsilon_0$. Maka $u_n
\to x_0$ tetapi $f(u_n) \not\to \ell$.
\end{proof}

\begin{corollary}[Operasi, komposisi, urutan]\label{cor:b1:continuity:operations}
Jumlah, hasil kali, dan hasil bagi (dengan limit penyebutnya taknol) atas limit berperilaku
seperti pada barisan; lalu jika $f \to \ell$ di $x_0$ dan $g \to m$ di $\ell$,
dan $g$ terdefinisi di sekitar $\ell$ dengan $g(\ell) = m$ atau $f \neq \ell$
di dekat $x_0$, maka $g \circ f \to m$ di $x_0$; adapun limit mengawetkan
ketaksamaan yang lebar, dan teorema apitnya berlaku.
\end{corollary}

\begin{proof}
Setiap \omterm{def:b1:logic:statement}{pernyataannya} berpindah lewat
\cref{thm:b1:continuity:seqchar} ke padanan barisannya
(\cref{thm:b1:seq:operations,thm:b1:seq:order}). Untuk komposisinya,
perangkaian langsungnya pantas ditulis sekali: diberikan $\varepsilon >
0$, limit $g$ di $\ell$ memasok $\eta > 0$ dengan
\[
\abs{y - \ell} \leq \eta \implies \abs{g(y) - m} \leq
\varepsilon \quad (y \text{ pada daerah asal } g),
\]
dengan kasus $y = \ell$ tercakup karena $g(\ell) = m$; lalu
limit $f$ di $x_0$ memasok $\delta > 0$ dengan $\abs{x -
x_0} \leq \delta \implies \abs{f(x) - \ell} \leq \eta$; sehingga dirangkai,
$\abs{x - x_0} \leq \delta$ memberikan $\abs{g(f(x)) - m} \leq
\varepsilon$. Adapun pada hipotesis alternatifnya ($f \neq \ell$ di dekat
$x_0$), nilai $y = \ell$ tak pernah disuapkan ke $g$ sehingga nilainya
di sana tak relevan.
\end{proof}

\begin{example}[Mengapa syarat pada komposisinya ada]\label{ex:b1:continuity:compositiontrap}
Misalkan $g(y) = 0$ untuk $y \neq 0$ dan $g(0) = 1$, dan $f$
identik dengan $0$. Maka $f(x) \to 0$ ketika $x \to 0$, dan $g(y) \to 0$
ketika $y \to 0$; namun $g(f(x)) = g(0) = 1$ untuk setiap $x$, sehingga $g
\circ f \to 1 \neq 0$. Fungsi dalamnya duduk \emph{persis pada}
nilai terlarang $\ell = 0$ selamanya, dan limit $g$ di $0$
mengabaikan apa yang dilakukan $g$ di $0$. Adapun syarat pada
\cref{cor:b1:continuity:operations} --- entah $g(\ell) = m$
(yakni $g$ \omterm{def:b1:continuity:continuous}{kontinu} di $\ell$), entah $f \neq \ell$ di dekat $x_0$ ---
persis itulah yang menyingkirkannya. Inti gagasan penutupnya: bahwa dalam
praktik orang menyusun fungsi yang \emph{\omterm{def:b1:continuity:continuous}{kontinu}} sehingga
syaratnya cuma-cuma; ia menggigit hanya ketika limitnya diambil sepanjang
\omterm{def:b1:topology:open}{persekitaran} yang berlubang, dan itulah sebabnya definisi
$\lim_{x \to x_0}$ yang dipakai buku ini menyertakan titiknya bila ia
termasuk daerah asalnya.
\end{example}

\section{Kekontinuan}

\begin{definition}\label{def:b1:continuity:continuous}
Fungsi $f$ disebut \emph{kontinu di $x_0 \in I$}\index{kekontinuan} bila $f(x)
\to f(x_0)$ ketika $x \to x_0$; dan \emph{kontinu pada $I$} bila ia
kontinu di setiap titiknya. Menurut
\cref{thm:b1:continuity:seqchar}: $f$ kontinu di $x_0$ jika dan hanya jika
$f(u_n) \to f(x_0)$ untuk setiap barisan $u_n \to x_0$ di $I$.
\end{definition}

\begin{example}[Taksonomi ketakkontinuan]\label{ex:b1:continuity:taxonomy}
Ada tiga cara gagal di sebuah titik, dengan keparahan yang menaik.
\emph{Terhapuskan}: karena $f(x) = \frac{\sin x}{x}$ pada $\R^*$ berlimit
$1$ di $0$; sehingga menetapkan $f(0) = 1$ memperbaikinya --- jadi ketakkontinuannya
tadi sebuah lubang, bukan ciri. \emph{Lompatan}: $\lfloor x \rfloor$ di
sebuah bilangan bulat mempunyai limit sepihak yang berbeda ($n - 1$ dan $n$); sehingga tak ada
pemilihan nilai yang dapat merukunkannya, tetapi kedua limit separuhnya ada.
\emph{Hakiki}: $\sin\frac1x$ di $0$ sama sekali tak mempunyai limit sepihak
(\cref{exo:b1:continuity:1}) --- yaitu ayunan tanpa
pengendapan. Adapun fungsi monoton hanya dapat menghasilkan jenis yang di tengah
(karena limit sepihaknya selalu ada, sebagai \omterm{def:b1:reals:bounds}{supremum} dan \omterm{def:b1:reals:bounds}{infimum}),
dan itulah sebabnya \omterm{def:b1:logic:sets}{himpunan} ketakkontinuannya paling banyak terbilang ---
yaitu satu bilangan rasional per lompatan. Sedangkan turunan, menurut teorema Darboux
(\cref{exo:b1:derivative:10}), hanya dapat menghasilkan jenis yang terakhir:
jadi fungsi dengan ketakkontinuan lompatan tak pernah menjadi turunan
apa pun.
\end{example}

\begin{proposition}\label{prop:b1:continuity:algebra}
Jumlah, hasil kali, hasil bagi (di tempat ia terdefinisi) dan komposisi
fungsi yang \omterm{def:b1:continuity:continuous}{kontinu} bersifat \omterm{def:b1:continuity:continuous}{kontinu}. Adapun \omterm{def:b1:poly:def}{polinomial}, \omterm{def:b1:fractions:field}{pecahan rasional}
(di luar polnya), $\abs{\,\cdot\,}$, $\exp$, $\ln$, fungsi
trigonometri dan \omterm{def:b1:functions:hyperbolic}{fungsi hiperbolik} beserta inversnya
(\cref{ch:b1:functions}) bersifat \omterm{def:b1:continuity:continuous}{kontinu} pada daerah asalnya.
\end{proposition}

\begin{proof}
Operasinya: \cref{cor:b1:continuity:operations}. Adapun konstanta dan
identitas bersifat \omterm{def:b1:continuity:continuous}{kontinu} langsung dari definisinya
($\delta = \varepsilon$ melayani identitasnya, dan sebarang $\delta$ bagi
konstantanya); lalu karena hasil kali fungsi yang \omterm{def:b1:continuity:continuous}{kontinu} bersifat \omterm{def:b1:continuity:continuous}{kontinu},
setiap monomial $a_k x^k$ menyusul lewat induksi pada $k$, dan jumlahnya
menuntaskan \omterm{def:b1:poly:def}{polinomialnya}; sedangkan \omterm{def:b1:fractions:field}{pecahan rasional} adalah hasil bagi dua
\omterm{def:b1:poly:def}{polinomial}, yang \omterm{def:b1:continuity:continuous}{kontinu} di mana pun penyebutnya tak
lenyap. Untuk $\abs{\,\cdot\,}$: lewat ketaksamaan segitiga terbalik,
$\bigl|\abs{f(x)} - \abs{f(x_0)}\bigr| \leq \abs{f(x) -
f(x_0)}$, sehingga $\delta$ yang sama berlaku. Adapun untuk fungsi
klasiknya kita menerima \omterm{def:b1:continuity:continuous}{kekontinuannya} di sini; sedangkan sifat dapat diturunkannya (yang dibuktikan pada
\cref{ch:b1:derivative}) lebih kuat.
\end{proof}

\begin{example}[Maks dan min fungsi yang kontinu]\label{ex:b1:continuity:maxmin}
Jika $f, g$ \omterm{def:b1:continuity:continuous}{kontinu}, maka $\max(f, g)$ dan $\min(f, g)$ pun \omterm{def:b1:continuity:continuous}{kontinu}:
jadi tanpa analisis kasus, berkat kesamaan
\[
\max(f, g) = \frac{f + g + \abs{f - g}}{2},
\qquad
\min(f, g) = \frac{f + g - \abs{f - g}}{2},
\]
dan \omterm{def:b1:continuity:continuous}{kekontinuan} jumlah beserta $\abs{\,\cdot\,}$
(\cref{prop:b1:continuity:algebra}). Khususnya $f^+ =
\max(f, 0)$ dan $f^- = \max(-f, 0)$ bersifat \omterm{def:b1:continuity:continuous}{kontinu} dengan $f =
f^+ - f^-$: jadi pemisahan tanda yang dipakai bagi deret
(\cref{ch:b1:series}) dan, dalam skala penuh, pada teori
pengintegralan jilid Tahun ke-3, tidak berongkos apa pun dalam keteraturannya.
\end{example}

\begin{theorem}[Teorema nilai antara]\label{thm:b1:continuity:ivt}
Misalkan $f$ \omterm{def:b1:continuity:continuous}{kontinu} pada $\intcc{a}{b}$ dengan $f(a) \leq 0 \leq
f(b)$. Maka $f(c) = 0$ untuk suatu $c \in \intcc{a}{b}$.
Akibatnya, fungsi yang \omterm{def:b1:continuity:continuous}{kontinu} pada sebuah \omterm{prop:b1:reals:intervals}{selang} mengambil setiap nilai
di antara sebarang dua nilainya: jadi $f(I)$ merupakan
\omterm{prop:b1:reals:intervals}{selang}.\index{teorema nilai antara}
\end{theorem}

\begin{proof}
Lewat dikotomi. Tetapkan $a_0 = a$, $b_0 = b$. Diberikan $\intcc{a_n}{b_n}$ dengan
$f(a_n) \leq 0 \leq f(b_n)$, misalkan $m$ titik tengahnya: jika $f(m) \leq
0$ pertahankan $\intcc{m}{b_n}$, dan kalau tidak pertahankan $\intcc{a_n}{m}$; maka syarat
tandanya bertahan. Barisan $(a_n), (b_n)$ itu berdampingan (karena $b_n -
a_n = \frac{b-a}{2^n}$), dengan limit bersama $c$
(\cref{thm:b1:seq:adjacent}). Lalu menurut \omterm{def:b1:continuity:continuous}{kekontinuan} dan
\cref{thm:b1:seq:order}: $f(c) = \lim f(a_n) \leq 0$ dan $f(c) =
\lim f(b_n) \geq 0$, sehingga $f(c) = 0$.

Untuk akibatnya: diberikan nilai $f(u) < v < f(w)$, terapkanlah yang di atas
pada $x \mapsto f(x) - v$ pada ruas yang bertitik ujung $u$ dan $w$;
sehingga $f(I)$ cembung, yakni sebuah \omterm{prop:b1:reals:intervals}{selang}
(\cref{prop:b1:reals:intervals}).
\end{proof}

\begin{omfigure}
\begin{tikzpicture}
\begin{axis}[omaxis, width=8.4cm, height=5.6cm,
    xmin=0, xmax=4.4, ymin=-1.1, ymax=1.4,
    xtick={0.6,3.8}, xticklabels={$a$,$b$}, ytick=\empty]
  \addplot[omDef, very thick, domain=0.6:3.8, samples=160]
    {0.5*(x-2) + 0.4*sin(deg(2.5*x))};
  \draw[dashed, gray] (axis cs:0.6,0) -- (axis cs:0.6,-0.301);
  \draw[dashed, gray] (axis cs:3.8,0) -- (axis cs:3.8,0.870);
  \node[below, font=\small, omDef] at (axis cs:0.6,-0.32) {$f(a)$};
  \node[above, font=\small, omDef] at (axis cs:3.8,0.89) {$f(b)$};
\end{axis}
\end{tikzpicture}

{\small Fungsi yang \omterm{def:b1:continuity:continuous}{kontinu} dengan $f(a) < 0 < f(b)$ pasti memotong
sumbunya: karena bukti dikotomi \cref{thm:b1:continuity:ivt} memerangkap sebuah
titik potong di antara barisan yang berdampingan.}
\end{omfigure}

\begin{example}[Satu persamaan, dengan protokol lengkapnya]\label{ex:b1:continuity:protocol}
Pecahkanlah $\eu^x = 3 - x$ di atas $\R$: yaitu keberadaan, ketunggalan, dan
letaknya. Tetapkan $g(x) = \eu^x + x - 3$, yang \omterm{def:b1:continuity:continuous}{kontinu}. Untuk letak dan
keberadaannya: $g(0) = -2 < 0$ dan $g(1) = \eu - 2 > 0$, sehingga
teorema nilai antara menanamkan sebuah penyelesaian di
$\intoo{0}{1}$. Ketunggalannya: karena $g$ merupakan jumlah $\eu^x$ yang
naik tegas dengan $x - 3$, sehingga ia naik tegas pada
$\R$; dan fungsi yang monoton tegas mengambil setiap nilainya paling banyak
sekali, sehingga penyelesaiannya tunggal pada seluruh $\R$ (bukan hanya pada
\omterm{prop:b1:reals:intervals}{selang} yang diperiksa). Adapun protokolnya --- susun ulang menjadi $g = 0$, lalu perubahan
tanda untuk keberadaannya, lalu kemonotonan untuk ketunggalannya --- menyelesaikan
kebanyakan pertanyaan ``berapa banyak penyelesaiannya'' dalam tiga baris, dan
tabel variasi pada \cref{ch:b1:derivative}
memperluasnya ke $g$ yang tak monoton dengan memotong $\R$ menjadi cabang
yang monoton.
\end{example}

\begin{example}[Dikotomi sebagai algoritma]\label{ex:b1:continuity:dichotomy}
Bukti \cref{thm:b1:continuity:ivt} itu menghitung. Ambillah $f(x) =
x^3 + x - 1$: di sini $f(0) = -1 < 0 < 1 = f(1)$, sehingga sebuah akar terletak di
$\intoo{0}{1}$. Lalu memaruhkannya:
\[
f(0.5) = -0.375 < 0, \qquad
f(0.75) = 0.171875 > 0, \qquad
f(0.625) = -0.130859375 < 0 ,
\]
sehingga akarnya berturut-turut terperangkap di $\intoo{0.5}{1}$, lalu
$\intoo{0.5}{0.75}$, lalu $\intoo{0.625}{0.75}$ (dengan nilai sejatinya
$c \approx 0.6823$). Setelah $n$ langkah galatnya paling banyak
$\frac{b - a}{2^n}$: jadi sepuluh langkah memberikan tiga desimal, dan dua puluh memberikan
enam. Inti gagasan penutupnya: bahwa teorema nilai antara bukan
hanya \omterm{def:b1:logic:statement}{pernyataan} keberadaan --- karena bukti dikotominya merupakan
algoritma pencarian akar yang terjamin, meski lamban, yang terhadapnya
metode Newton pada \cref{ch:b1:derivative} yang cepat tetapi lokal patut
diukur.
\end{example}

\begin{theorem}[Teorema nilai ekstrem]\label{thm:b1:continuity:evt}
Fungsi yang \omterm{def:b1:continuity:continuous}{kontinu} pada sebuah ruas $\intcc{a}{b}$ bersifat terbatas dan
mencapai batasnya: yakni ada $c, d \in \intcc{a}{b}$ dengan
\[
f(c) = \inf_{\intcc{a}{b}} f,
\qquad
f(d) = \sup_{\intcc{a}{b}} f .
\]
Digabung dengan \cref{thm:b1:continuity:ivt}: \emph{peta ruas
di bawah fungsi yang \omterm{def:b1:continuity:continuous}{kontinu} adalah ruas} $\intcc{f(c)}{f(d)}$.
\index{teorema nilai ekstrem}
\end{theorem}

\begin{proof}
\emph{Terbatas di atas:} kalau tidak, pilihlah $u_n$ dengan $f(u_n) \geq n$. Menurut
kekompakan ruasnya (\cref{thm:b1:topology:compact}), suatu
\omterm{def:b1:seq:subsequence}{subbarisan} $u_{\varphi(n)} \to x \in \intcc{a}{b}$; lalu \omterm{def:b1:continuity:continuous}{kekontinuannya} memberikan
$f(u_{\varphi(n)}) \to f(x)$, padahal $f(u_{\varphi(n)}) \geq \varphi(n)
\to +\infty$: yang bertentangan.

\emph{Sup tercapai:} misalkan $M = \sup f$ lalu pilihlah $v_n$ dengan $f(v_n) >
M - \frac{1}{n+1}$ (\cref{prop:b1:reals:epsilon}). Ekstraklah
$v_{\varphi(n)} \to d \in \intcc{a}{b}$: maka $f(d) = \lim
f(v_{\varphi(n)}) = M$ menurut apitannya. Adapun \omterm{def:b1:reals:bounds}{infimumnya} ditangani oleh
$-f$.
\end{proof}

\begin{example}[Minimum yang positif pada sebuah ruas]\label{ex:b1:continuity:posmin}
Misalkan $f$ \omterm{def:b1:continuity:continuous}{kontinu} pada $\intcc{0}{1}$ dengan $f(x) > 0$ untuk
setiap $x$. Maka $\inf f > 0$: karena menurut teorema nilai ekstrem
\omterm{def:b1:reals:bounds}{infimumnya} merupakan sebuah \emph{nilai} $f(c)$, dan $f(c) > 0$ menurut hipotesisnya.
Jadi fungsi positif yang \omterm{def:b1:continuity:continuous}{kontinu} pada sebuah ruas terbatas menjauhi
$0$ --- yaitu argumen dua baris yang dipakai belasan kali pada
bab berikutnya (penyebut yang terkendali, pengapitan fungsi
tangga, batas galat). Pada \omterm{prop:b1:reals:intervals}{selang} yang tak kompak hal ini gagal
secara mencolok: karena $f(x) = x$ pada $\intoc{0}{1}$ \omterm{def:b1:continuity:continuous}{kontinu} dan
positif dengan $\inf f = 0$, yang tak tercapai. Inti gagasan penutupnya:
bahwa ``positif'' naik pangkat menjadi ``positif seragam'' tepat jika
daerah asalnya kompak; jadi setiap hipotesis teorema nilai ekstrem
memikul beban.
\end{example}

\begin{remark}[Jebakan yang lazim di sekitar ketiga teoremanya]
(i) \emph{Peta di bawah fungsi \omterm{def:b1:continuity:continuous}{kontinu}}: hanya \emph{ruas} yang kokoh.
Karena peta \omterm{prop:b1:reals:intervals}{selang} \omterm{def:b1:topology:open}{terbuka} di bawah fungsi yang \omterm{def:b1:continuity:continuous}{kontinu} tak harus \omterm{def:b1:topology:open}{terbuka}
($x \mapsto x^2$ memetakan $\intoo{-1}{1}$ pada $\intco{0}{1}$), dan
peta \omterm{def:b1:topology:closed}{himpunan tertutup} tak harus \omterm{def:b1:topology:closed}{tertutup} ($\arctan$ memetakan
$\R$ yang \omterm{def:b1:topology:closed}{tertutup} pada $\intoo{-\frac\pi2}{\frac\pi2}$ yang
\omterm{def:b1:topology:open}{terbuka}); tetapi peta sebuah ruas adalah
ruas (\cref{thm:b1:continuity:evt}). (ii) \emph{Teorema
nilai antara memerlukan sebuah \omterm{prop:b1:reals:intervals}{selang}}: karena fungsi $x
\mapsto \frac1x$, yang \omterm{def:b1:continuity:continuous}{kontinu} pada $\intco{-1}{0} \cup
\intoc{0}{1}$, mengambil nilai $-1$ dan $1$ namun tak pernah
lenyap --- sebab daerah asalnya dua kepingan yang lepas, dan nilai
$0$ jatuh ke celahnya; jadi namailah selalu \omterm{prop:b1:reals:intervals}{selang} yang di atasnya
teoremanya diterapkan. (iii) \emph{\omterm{def:b1:continuity:uniform}{Kekontinuan seragam} adalah sifat
pasangan (fungsi, \omterm{def:b1:logic:sets}{himpunan})}: karena $x^2$ \omterm{def:b1:continuity:uniform}{kontinu seragam} pada setiap
ruas namun tidak pada $\R$
(\cref{ex:b1:continuity:notuniform}) --- jadi kata ``\omterm{def:b1:continuity:uniform}{kontinu
seragam}'' tanpa daerah asalnya tak bermakna. (iv)
\emph{\omterm{def:b1:continuity:continuous}{Kekontinuan} inversnya tak bersifat formal}: ia berlaku pada
\omterm{prop:b1:reals:intervals}{selang} lewat kemonotonan
(\cref{thm:b1:continuity:bijection}), tetapi bijeksi yang \omterm{def:b1:continuity:continuous}{kontinu}
antara gabungan \omterm{prop:b1:reals:intervals}{selang} dapat berinvers takkontinu ---
adapun catatan setelah teorema itu ada karena mahasiswa mengutipnya
tanpa hipotesis \omterm{prop:b1:reals:intervals}{selangnya}.
\end{remark}

\section{Fungsi monoton dan fungsi invers}

\begin{theorem}[Teorema bijeksi monoton]\label{thm:b1:continuity:bijection}
Misalkan $f$ \omterm{def:b1:continuity:continuous}{kontinu} dan monoton tegas pada sebuah \omterm{prop:b1:reals:intervals}{selang} $I$. Maka
\begin{enumerate}
  \item $f$ merupakan bijeksi dari $I$ pada \omterm{prop:b1:reals:intervals}{selang} $J = f(I)$;
  \item inversnya $f^{-1} \colon J \to I$ bersifat monoton tegas (pada
        arah yang sama) dan \emph{\omterm{def:b1:continuity:continuous}{kontinu}}.
\end{enumerate}
\end{theorem}

\begin{proof}
Katakanlah $f$ naik tegas. (1) Keinjektifannya langsung dari
kemonotonan tegasnya; kesurjektifannya pada $f(I)$ sepele, dan $f(I)$
merupakan \omterm{prop:b1:reals:intervals}{selang} menurut \cref{thm:b1:continuity:ivt}.

(2) Adapun $f^{-1}$ naik tegas: karena jika $y < y'$ di $J$ tetapi $f^{-1}(y)
\geq f^{-1}(y')$, maka menerapkan $f$ yang naik memberikan $y \geq y'$,
yang mustahil. \omterm{def:b1:continuity:continuous}{Kekontinuan} $f^{-1}$ di $y_0 = f(x_0) \in J$: misalkan
$\varepsilon > 0$. Andaikan mula-mula $x_0$ \omterm{def:b1:topology:closure}{interior} bagi $I$, lalu
kecilkan $\varepsilon$ sehingga $x_0 \pm \varepsilon \in I$: maka
petanya memenuhi $f(x_0 - \varepsilon) < y_0 < f(x_0 +
\varepsilon)$. Ambillah
$\delta = \min\bigl(y_0 - f(x_0 - \varepsilon),\, f(x_0 +
\varepsilon) - y_0\bigr) > 0$: sehingga untuk $\abs{y - y_0} \leq \delta$,
kemonotonan $f^{-1}$ mengapit $f^{-1}(y)$ di antara $x_0 -
\varepsilon$ dan $x_0 + \varepsilon$. Jika $x_0$, katakanlah, titik ujung
kiri $I$, maka hanya $x_0 + \varepsilon$ yang tersedia: lalu
$y_0 = \min J$ (menurut kemonotonannya), setiap $y \in J$ dengan $y - y_0
\leq f(x_0 + \varepsilon) - y_0$ memenuhi $x_0 \leq f^{-1}(y)
\leq x_0 + \varepsilon$, dan taksiran sepihaknya persis merupakan
\omterm{def:b1:continuity:continuous}{kekontinuan} di sebuah titik ujung; adapun titik ujung kanannya simetris.
(Perhatikan: \omterm{def:b1:continuity:continuous}{kekontinuan} $f^{-1}$ \emph{tak} disimpulkan dari
\omterm{def:b1:continuity:continuous}{kekontinuan} $f$ lewat kesimetrian --- karena kemonotonan
pada sebuah selanglah yang mengerjakannya.)
\end{proof}

\begin{remark}
Teorema inilah yang dipakai diam-diam oleh \cref{def:b1:functions:arc} dan
\cref{prop:b1:functions:invhyp}: bahwa $\arcsin$, $\arccos$,
$\arctan$, $\operatorname{arsinh}$, \dots\ bersifat \omterm{def:b1:continuity:continuous}{kontinu}. Ada
pelengkapnya (\cref{exo:b1:continuity:10}): bahwa fungsi
\emph{\omterm{def:b1:logic:inj}{injektif}} yang \omterm{def:b1:continuity:continuous}{kontinu} pada sebuah \omterm{prop:b1:reals:intervals}{selang} otomatis monoton
tegas, sehingga hipotesis kemonotonannya tak berongkos apa pun.
\end{remark}

\begin{example}[Akar segala orde]\label{ex:b1:continuity:nthroot}
Untuk $n \in \N^*$, fungsi $f(x) = x^n$ bersifat \omterm{def:b1:continuity:continuous}{kontinu} dan
naik tegas pada $\intco{0}{+\infty}$, dengan $f(0) = 0$ dan
$f(x) \to +\infty$: sehingga petanya adalah $\intco{0}{+\infty}$
(menurut \cref{thm:b1:continuity:ivt} untuk struktur \omterm{prop:b1:reals:intervals}{selangnya}). Lalu
teorema bijeksi monotonnya menghasilkan, dalam satu tarikan, sebuah
invers \emph{\omterm{def:b1:continuity:continuous}{kontinu}} yang naik tegas
\[
x \mapsto x^{1/n} \colon \intco{0}{+\infty} \to
\intco{0}{+\infty} :
\]
yakni keberadaan, ketunggalan dan \omterm{def:b1:continuity:continuous}{kekontinuan} akar pangkat $n$, tanpa
perhitungan apa pun. Bandingkanlah dengan \cref{exo:b1:reals:12}, yang membangun
$\sqrt y$ dengan tangan dari \omterm{def:b1:reals:bounds}{supremumnya}: jadi satu bab teori sudah
memampatkan satu halaman kerja itu menjadi dua baris, dan dua baris yang
sama sudah mengesahkan $\arcsin$, $\arctan$ dan
$\operatorname{arsinh}$ sebelumnya. Inti gagasan penutupnya: bahwa teorema
yang baik adalah kerja yang tersimpan.
\end{example}

\section{Kekontinuan seragam}

\begin{definition}\label{def:b1:continuity:uniform}
Fungsi $f \colon I \to \R$ disebut \emph{kontinu seragam}\index{kekontinuan
seragam} bila
\[
\forall \varepsilon > 0,\ \exists \delta > 0,\ \forall x, y \in I,
\qquad \abs{x - y} \leq \delta \implies \abs{f(x) - f(y)} \leq
\varepsilon .
\]
Intinya: bahwa $\delta$ bergantung pada $\varepsilon$ saja, bukan pada
letaknya di $I$. \omterm{def:b1:continuity:continuous}{Kekontinuan} seragam mengakibatkan \omterm{def:b1:continuity:continuous}{kekontinuan}; dan
fungsi \emph{Lipschitz} ($\abs{f(x) - f(y)} \leq k \abs{x - y}$)
bersifat \omterm{def:b1:continuity:continuous}{kontinu} seragam (dengan $\delta = \varepsilon/k$).
\end{definition}

\begin{example}\label{ex:b1:continuity:notuniform}
Fungsi $x \mapsto x^2$ \omterm{def:b1:continuity:continuous}{kontinu} pada $\R$ tetapi \emph{tak} \omterm{def:b1:continuity:uniform}{kontinu
seragam}: karena $\abs{(n + \frac 1n)^2 - n^2} = 2 + \frac{1}{n^2} \geq 2$
padahal argumennya berjarak $\frac 1n \to 0$. Namun pada sebarang
\omterm{prop:b1:reals:intervals}{selang} yang terbatas ia Lipschitz, sehingga \omterm{def:b1:continuity:uniform}{kontinu seragam} ---
yang selaras dengan teorema Heine di bawah.
\end{example}

\begin{example}[Modulus keseragamannya, secara eksplisit]\label{ex:b1:continuity:moduli}
Pada sebuah ruas, Heine menjamin $\delta$ yang seragam; dan sering kali kita
juga dapat \emph{menghitungnya}. Untuk $f(x) = x^2$ pada $\intcc{0}{10}$:
\[
\abs{x^2 - y^2} = \abs{x + y}\,\abs{x - y} \leq 20\,\abs{x - y},
\]
sehingga $\delta = \frac{\varepsilon}{20}$ berlaku secara seragam (yaitu
\omterm{def:b1:complex:field}{modulus} Lipschitz, yang linear dalam $\varepsilon$). Untuk $\sqrt x$ pada
$\intcc{0}{1}$: $\abs{\sqrt x - \sqrt y} \leq \sqrt{\abs{x -
y}}$ (\cref{exo:b1:continuity:9}), sehingga $\delta = \varepsilon^2$
berlaku --- yang seragam tetapi \emph{tak} linear: karena di dekat $0$ akar
kuadratnya curam, dan harganya muncul pada eksponen
$\varepsilon$, bukan pada kegagalan keseragamannya. Inti gagasan
penutupnya: bahwa \omterm{def:b1:continuity:uniform}{kekontinuan seragam} adalah spektrum, bukan ya/tidak ---
karena fungsi $\delta(\varepsilon)$, yang disebut modulusnya, mengukur
seberapa mahal keseragaman itu, dan Lipschitz sekadar peringkatnya
yang terbaik.
\end{example}

\begin{theorem}[Heine]\label{thm:b1:continuity:heine}
Fungsi yang \omterm{def:b1:continuity:continuous}{kontinu} pada sebuah \emph{ruas} $\intcc{a}{b}$ bersifat \omterm{def:b1:continuity:uniform}{kontinu
seragam}.\index{teorema Heine}
\end{theorem}

\begin{proof}
Lewat pertentangan: andaikan suatu $\varepsilon_0 > 0$ mengalahkan setiap
$\delta$. Dengan $\delta = \frac{1}{n+1}$, pilihlah $x_n, y_n \in
\intcc{a}{b}$ dengan $\abs{x_n - y_n} \leq \frac{1}{n+1}$ dan
$\abs{f(x_n) - f(y_n)} > \varepsilon_0$. Lalu kekompakannya
(\cref{thm:b1:topology:compact}) mengekstrak $x_{\varphi(n)} \to c \in
\intcc{a}{b}$; sehingga $y_{\varphi(n)} \to c$ pula (lewat apitan pada $\abs{x -
y}$). \omterm{def:b1:continuity:continuous}{Kekontinuan} di $c$ memberikan $f(x_{\varphi(n)}) \to f(c)$ dan
$f(y_{\varphi(n)}) \to f(c)$, sehingga $\abs{f(x_{\varphi(n)}) -
f(y_{\varphi(n)})} \to 0 < \varepsilon_0$: yang bertentangan.
\end{proof}

\begin{example}[Terbatas, kontinu, namun tak seragam]\label{ex:b1:continuity:sinx2}
Fungsi $f(x) = \sin(x^2)$ bersifat \omterm{def:b1:continuity:continuous}{kontinu} dan terbatas pada
$\R$, tetapi tak \omterm{def:b1:continuity:uniform}{kontinu seragam}. Ambillah
\[
x_n = \sqrt{2\pi n}, \qquad y_n = \sqrt{2\pi n + \tfrac\pi2}:
\qquad
y_n - x_n = \frac{\pi/2}{\sqrt{2\pi n + \frac\pi2} +
\sqrt{2\pi n}} \longrightarrow 0 ,
\]
namun $f(y_n) - f(x_n) = \sin\bigl(2\pi n + \frac\pi2\bigr) -
\sin(2\pi n) = 1 - 0 = 1$ untuk setiap $n$: sehingga tak ada $\delta$ tunggal yang dapat
melayani $\varepsilon = \frac12$ di mana-mana. Secara geometri,
ayunan $\sin(x^2)$ \emph{mempercepat}: karena grafiknya
menuntaskan satu gelombang penuh pada jendela yang kian pendek, sehingga
skala mendatar yang dituntut sebuah $\varepsilon$ menyusut ke nol
ketika $x$ tumbuh. Inti gagasan penutupnya: bahwa keterbatasan tak membeli
keseragaman (menurut contoh ini), dan ketakterbatasan tak menghalanginya
($\sqrt x$, \cref{exo:b1:continuity:9}); adapun yang memutuskan adalah
\emph{\omterm{def:b1:complex:field}{modulus} ayunannya}, dan teorema Heine mengatakan bahwa daerah asal yang
kompak mendisiplinkannya secara otomatis.
\end{example}

\begin{example}[Percobaan kalkulator, yang dijelaskan]\label{ex:b1:continuity:cositeration}
Ketiklah sebarang bilangan pada kalkulator lalu tekanlah $\cos$ berulang-ulang:
maka layarnya mengunci pada $0.7390851\dots$ Mengapa? Setelah satu tekanan
nilainya terletak di $\intcc{-1}{1}$, setelah dua di $\intcc{\cos 1}{1}
\subseteq \intcc{0.54}{1}$, yaitu \omterm{prop:b1:reals:intervals}{selang} yang stabil bagi $\cos$. Di atasnya,
$\abs{\cos a - \cos b} \leq \sin(1)\,\abs{a - b}$ dengan $\sin 1 =
0.841\dots < 1$ (menurut batas hasil kali menjadi jumlah pada
\cref{pb:b1:seq:1}, atau ketaksamaan nilai rata-rata pada
\cref{ch:b1:derivative}): sehingga iterasinya mengontraksi, jadi menurut
langkah kendali galat pada \cref{met:b1:seq:recurrent},
\[
\abs{u_n - c} \leq (0.842)^{\,n-2}\,\abs{u_2 - c}
\longrightarrow 0 ,
\]
dengan $c$ titik tetap tunggal $\cos c = c$
(\cref{exo:b1:continuity:6}). Kira-kira $40$ tekanan membeli tiga
desimal (karena $0.842^{40} \approx 10^{-3}$) --- yaitu laju geometri,
yang lebih jinak daripada $2^{-n}$ per langkah pada dikotominya, tetapi setiap tekanannya berongkos
satu ketukan tombol sedangkan setiap langkah dikotominya berongkos satu
penilaian tanda yang penuh. Inti gagasan penutupnya: bahwa gambaran titik tetap pada
\cref{ch:b1:seq} dan teorema keberadaan pada bab ini merupakan
dua paruh dari satu kisah --- teorema nilai antara menemukan $c$, sedangkan kontraksinya mencapainya.
\end{example}

\begin{remark}[Di mana teorema ini bekerja berikutnya]
Ketiga pilar bab ini masing-masing menggerakkan bab yang belakangan. Adapun
teorema nilai antara menyuapi setiap argumen keberadaan penyelesaian
dan teorema bijeksi monotonnya; sedangkan teorema nilai
ekstrem mengubah masalah pengoptimuman menjadi teorema (karena Rolle dan
teorema nilai rata-rata pada \cref{ch:b1:derivative} bermula persis
di situ); sementara teorema Heine adalah alasan mengapa fungsi yang \omterm{def:b1:continuity:continuous}{kontinu} pada
ruas dapat diintegralkan pada \cref{ch:b1:integration} --- karena
$\delta$ yang seragam itulah yang membuat jumlah Riemannnya konvergen. Pada jilid Tahun
ke-2 trio yang sama muncul kembali pada ruang vektor bernorma, dengan
kekompakan mengerjakan apa yang dikerjakan ruas di sini.
\end{remark}

\begin{remark}[Cakrawala di dalam jilid ini]
\omterm{def:b1:continuity:continuous}{Kekontinuan} segera akan dilampaui tetapi tak pernah dipensiunkan.
\Cref{ch:b1:derivative} memperkuatnya menjadi sifat dapat diturunkan lalu
membalas jasanya (karena dapat diturunkan mengakibatkan \omterm{def:b1:continuity:continuous}{kontinu}); sedangkan
\cref{ch:b1:integration} bersandar padanya dua kali, lewat Heine untuk
konstruksinya dan lewat teorema fundamentalnya, yang objek
pusatnya $x \mapsto \int_a^x f$ menaikkan pangkat $f$ yang sekadar
\omterm{def:b1:continuity:continuous}{kontinu} menjadi primitif $C^1$. Pada
\cref{ch:b1:multivar}, \omterm{def:b1:continuity:continuous}{kekontinuan} dalam dua variabel menyimpan jebakan
yang pantas ditengok lebih dulu: bahwa fungsi $\frac{xy}{x^2 + y^2}$ (yang diperluas
oleh $0$) bersifat \omterm{def:b1:continuity:continuous}{kontinu} dalam $x$ untuk setiap $y$ yang tetap dan dalam $y$ untuk
setiap $x$ yang tetap, namun tak \omterm{def:b1:continuity:continuous}{kontinu} di titik asalnya --- karena sepanjang
diagonal $x = y$ ia selalu $\frac12$. Jadi \omterm{def:b1:continuity:continuous}{kekontinuan} yang terpisah
sungguh lebih lemah daripada \omterm{def:b1:continuity:continuous}{kekontinuan}: pencirian lewat
barisan memang bertahan pada perpindahan ke $\R^2$, tetapi barisannya harus
diizinkan mendekat dari \emph{setiap} arah, bukan hanya
sepanjang sumbunya.
\end{remark}

\section{Latihan}

\begin{exercise}[$\star$]\label{exo:b1:continuity:1}
Dengan memakai pencirian lewat barisan, buktikan bahwa $x \mapsto
\sin\frac 1x$ tak berlimit di $0^+$ \emph{(tunjukkanlah dua barisan)}.
Apakah $x \mapsto x \sin\frac 1x$ mempunyai limit?
\end{exercise}

\begin{exercise}[$\star$]\label{exo:b1:continuity:2}
Telaahlah \omterm{def:b1:continuity:continuous}{kekontinuan} pada $\R$ bagi $f(x) = \lfloor x \rfloor$ dan bagi
$g(x) = x - \lfloor x \rfloor$, dan bagi $h(x) = \lfloor x \rfloor +
(x - \lfloor x\rfloor)^2$.
\end{exercise}

\begin{exercise}[$\star$]\label{exo:b1:continuity:3}
Buktikan bahwa persamaan $x^5 - 3x + 1 = 0$ mempunyai sekurang-kurangnya tiga penyelesaian
real \emph{(nilailah pada titik yang terpilih baik lalu terapkan
\cref{thm:b1:continuity:ivt} pada tiga ruas yang lepas)}.
\end{exercise}

\begin{exercise}[$\star$]\label{exo:b1:continuity:4}
Buktikan bahwa setiap \omterm{def:b1:poly:def}{polinomial} berderajat ganjil mempunyai akar real.
\end{exercise}

\begin{exercise}[$\star\star$]\label{exo:b1:continuity:5}
(Titik tetap) Misalkan $f \colon \intcc{0}{1} \to \intcc{0}{1}$
\omterm{def:b1:continuity:continuous}{kontinu}. Buktikan bahwa $f$ mempunyai titik tetap: yakni $f(c) = c$ untuk suatu
$c$. Perlihatkanlah bahwa \omterm{def:b1:continuity:continuous}{kekontinuannya} maupun ruasnya tak dapat
dilepaskan.
\end{exercise}

\begin{exercise}[$\star\star$]\label{exo:b1:continuity:6}
Buktikan bahwa persamaan $\cos x = x$ mempunyai tepat satu penyelesaian real,
dan bahwa ia terletak di $\intoo{0}{1}$.
\end{exercise}

\begin{exercise}[$\star\star$]\label{exo:b1:continuity:7}
Misalkan $f \colon \R \to \R$ \omterm{def:b1:continuity:continuous}{kontinu} dengan $f(x) \to +\infty$ ketika $x
\to \pm\infty$. Buktikan bahwa $f$ mencapai minimum semestanya pada $\R$.
\emph{(Reduksikanlah ke sebuah ruas yang memuat \omterm{def:b1:logic:sets}{himpunan} subarasnya.)}
\end{exercise}

\begin{exercise}[$\star\star$]\label{exo:b1:continuity:8}
Misalkan $f \colon \R \to \R$ \omterm{def:b1:continuity:continuous}{kontinu} dan periodik (dengan periode $T >
0$). Buktikan bahwa $f$ terbatas dan mencapai batasnya, dan bahwa
ada $c$ dengan $f(c + \frac T2) = f(c)$. \emph{(Untuk butir
keduanya, telaahlah $g(x) = f(x + \frac T2) - f(x)$ pada satu periode.)}
\end{exercise}

\begin{exercise}[$\star\star$]\label{exo:b1:continuity:9}
Buktikan bahwa $x \mapsto \sqrt x$ \omterm{def:b1:continuity:uniform}{kontinu seragam} pada
$\intco{0}{+\infty}$, meskipun ia tak Lipschitz di dekat $0$.
\emph{(Buktikan lalu pakailah $\abs{\sqrt x - \sqrt y} \leq \sqrt{\abs{x -
y}}$.)}
\end{exercise}

\begin{exercise}[$\star\star\star$]\label{exo:b1:continuity:10}
Misalkan $f$ \omterm{def:b1:continuity:continuous}{kontinu} dan \omterm{def:b1:logic:inj}{injektif} pada sebuah \omterm{prop:b1:reals:intervals}{selang} $I$. Buktikan bahwa
$f$ monoton tegas. \emph{Petunjuk: kalau tidak, ada $a < b < c$
dengan, katakanlah, $f(b) > f(a)$ dan $f(b) > f(c)$; lalu terapkan teorema
nilai antara pada sebuah nilai di antara $\max(f(a), f(c))$ dan $f(b)$ pada
kedua sisi $b$.}
\end{exercise}

\begin{exercise}[$\star\star\star$]\label{exo:b1:continuity:11}
(Persamaan fungsional Cauchy, kasus yang \omterm{def:b1:continuity:continuous}{kontinu}) Misalkan $f \colon \R \to
\R$ \omterm{def:b1:continuity:continuous}{kontinu} dengan $f(x + y) = f(x) + f(y)$ untuk setiap $x, y$.
Buktikan bahwa $f(x) = f(1)\,x$ untuk setiap $x$: mula-mula pada $\N$, $\Z$, $\Q$
(lewat keaditifannya belaka), lalu pada $\R$ lewat \omterm{def:b1:continuity:continuous}{kekontinuan} dan \omterm{def:b1:topology:dense}{kepadatan}
(\cref{thm:b1:reals:density}).
\end{exercise}

\begin{exercise}[$\star\star\star$]\label{exo:b1:continuity:12}
Misalkan $f \colon \intco{0}{+\infty} \to \R$ \omterm{def:b1:continuity:continuous}{kontinu} dengan
$f(x) \to \ell \in \R$ ketika $x \to +\infty$. Buktikan bahwa $f$
\omterm{def:b1:continuity:continuous}{kontinu} \emph{seragam} pada $\intco{0}{+\infty}$. \emph{(Potonglah pada $A$
yang besar: Heine pada $\intcc{0}{A+1}$, dan limitnya di luar $A$; lalu buatlah
kedua rezimnya bertindihan.)}
\end{exercise}

\section{Soal: Persamaan fungsional Cauchy dan saudarinya}

\begin{problem}[{Soal akhir pekan --- $f(x + y) = f(x) + f(y)$:
keteraturan memaksa kelinearan, beserta potret monsternya}]
\label{pb:b1:continuity:1}
Fungsi mana yang memenuhi $f(x + y) = f(x) + f(y)$ untuk setiap $x,
y$ yang real? Cauchy mengajukan pertanyaannya pada tahun 1821; dan jawabannya sebuah paradigma.
\Cref{exo:b1:continuity:11} menunjukkan bahwa fungsi \emph{aditif}
yang demikian bersifat linear pada $\Q$ dan bahwa \omterm{def:b1:continuity:continuous}{kekontinuan} yang penuh memaksa $f(x)
= cx$. Soal ini mempertajam hipotesisnya secara dramatis ---
\omterm{def:b1:continuity:continuous}{kekontinuan} di \emph{satu} titik, atau kemonotonan, atau sekadar
keterbatasan pada satu \omterm{prop:b1:reals:intervals}{selang} yang kecil, masing-masingnya mencukupi --- lalu melukis
potret sebuah penyelesaian taklinear yang hipotetis (yang grafiknya
mengisi bidangnya), memecahkan persamaan saudarinya yang mencirikan
$\eu^{cx}$, $c\ln x$, $x^c$ dan $cx^2$, lalu menutup dengan persamaan Jensen
beserta teorema \emph{cembung titik tengah $+$ \omterm{def:b1:continuity:continuous}{kontinu}
$\implies$ cembung}.\index{persamaan fungsional
Cauchy}\index{persamaan fungsional} Di sepanjang soal ini, \emph{aditif}
berarti: $f(x + y) = f(x) + f(y)$ untuk setiap $x, y \in \R$.

\medskip
\noindent\textbf{Bagian I --- Kelinearan atas $\Q$, dan satu titik
\omterm{def:b1:continuity:continuous}{kekontinuan}.}
\begin{enumerate}
  \item Misalkan $f$ aditif. Menurut \cref{exo:b1:continuity:11},
        $f(r) = f(1)\,r$ untuk $r$ yang rasional. Buktikan \omterm{def:b1:logic:statement}{pernyataan}
        yang lebih halus yang dipakai di bawah: bahwa untuk \emph{setiap} $x \in \R$ dan $r
        \in \Q$, $f(rx) = r\,f(x)$ (jadi $f$ bersifat
        \emph{linear} atas $\Q$).
  \item Andaikan $f$ yang aditif itu \omterm{def:b1:continuity:continuous}{kontinu} di satu titik
        tunggal $x_0$. Tunjukkan bahwa $f$ \omterm{def:b1:continuity:continuous}{kontinu} di mana-mana
        \emph{(hitunglah $f(x + h) - f(x)$ dalam $f(x_0 + h)
        - f(x_0)$)}, sehingga $f(x) = f(1)\,x$.
  \item Tunjukkan bahwa fungsi yang aditif ditentukan oleh
        pembatasannya pada sebarang \omterm{def:b1:structures:subgroup}{subgrup} yang \omterm{def:b1:topology:dense}{padat}: bahwa jika dua fungsi
        yang aditif bersesuaian pada $\Z + \sqrt2\,\Z$ dan keduanya
        \omterm{def:b1:continuity:continuous}{kontinu}, maka keduanya sama --- sedangkan tanpa \omterm{def:b1:continuity:continuous}{kekontinuan},
        menetapkan $f(1) = 0$ dan $f(\sqrt 2) = 1$ tetap taat asas
        dengan kelinearan atas $\Q$ pada \omterm{def:b1:structures:subgroup}{subgrupnya}. Hitunglah $f(m +
        n\sqrt2)$ untuk penetapan ini.
  \item Misalkan $f$ aditif dan terbatas di atas oleh $M$ pada suatu
        \omterm{prop:b1:reals:intervals}{selang} $\intcc{a}{b}$ dengan $a < b$. Tunjukkan bahwa $f$
        terbatas di atas pada $\intcc{0}{\ell}$, dengan $\ell = b - a$
        \emph{(translasikanlah sejauh $a$)}.
\end{enumerate}

\medskip
\noindent\textbf{Bagian II --- Tangga keteraturannya.}
\begin{enumerate}[resume]
  \item Melanjutkan pertanyaan 4: dengan memakai $f(\ell - t) + f(t) =
        f(\ell)$, tunjukkan bahwa $f$ juga terbatas \emph{di bawah} pada
        $\intcc{0}{\ell}$: yakni $\abs f \leq C$ di sana.
  \item Tunjukkan bahwa $\abs{f(t)} \leq \frac{C}{n}$ untuk $t \in
        \intcc{0}{\ell/n}$, lalu simpulkan bahwa $f$ \omterm{def:b1:continuity:continuous}{kontinu} di
        $0$ \emph{(keganjilannya menangani sisi kirinya)}, sehingga
        di mana-mana (pertanyaan 2): jadi fungsi aditif yang terbatas pada
        satu \omterm{prop:b1:reals:intervals}{selang} bersifat linear.
  \item Simpulkan kasus monotonnya: bahwa fungsi aditif yang
        tidak turun pada suatu $\intcc{a}{b}$ ($a < b$) berbentuk $f(x)
        = cx$ dengan $c \geq 0$.
  \item Rakitlah \emph{tangga keteraturannya}: bahwa untuk $f$ yang
        aditif, keempat ini setara --- (a) $f(x) = cx$;
        (b) $f$ \omterm{def:b1:continuity:continuous}{kontinu}; (c) $f$ \omterm{def:b1:continuity:continuous}{kontinu} di satu titik; (d)
        $f$ monoton pada suatu \omterm{prop:b1:reals:intervals}{selang} yang taktermerosot; (e) $f$
        terbatas pada suatu \omterm{prop:b1:reals:intervals}{selang} yang taktermerosot. Aturlah
        implikasinya sehingga masing-masingnya entah sepele entah sudah
        terbukti.
  \item Periksalah bahwa pertanyaan 4--6 hanya memakai batas
        \emph{atasnya}: jadi fungsi aditif yang terbatas di atas pada satu
        \omterm{prop:b1:reals:intervals}{selang} taktermerosot sudah linear. Simpulkanlah
        \omterm{def:b1:logic:statement}{pernyataan} cerminnya bagi batas bawahnya, lalu catatlah
        bentuk terkuat anak tangga (e) yang diperoleh dengan demikian.
\end{enumerate}

\medskip
\noindent\textbf{Bagian III --- Potret sebuah monster.} Andaikan
sekarang $f$ aditif tetapi \emph{tak} linear.
\begin{enumerate}[resume]
  \item Tunjukkan bahwa ada bilangan real taknol $u, v$ dengan
        $\dfrac{f(u)}{u} \neq \dfrac{f(v)}{v}$, dan bahwa
        vektor $(u, f(u))$ dan $(v, f(v))$ merentang bidangnya (karena
        determinannya $u f(v) - v f(u)$ taknol).
  \item Tunjukkan bahwa grafik $f$ memuat semua titik
        \[
        r\,(u, f(u)) + s\,(v, f(v)), \qquad r, s \in \Q ,
        \]
        lalu simpulkan bahwa grafiknya \emph{\omterm{def:b1:topology:dense}{padat} di $\R^2$}: yakni untuk
        setiap titik $(x_0, y_0)$ pada bidangnya dan setiap
        $\varepsilon > 0$, ada $(x, f(x))$ yang berjarak
        $\varepsilon$ darinya \emph{(pecahkanlah sistem real $2\times2$-nya,
        lalu hampiri koefisien realnya dengan
        bilangan rasional)}.
  \item Simpulkanlah dari pertanyaan 11 potretnya yang lengkap: bahwa fungsi
        aditif yang taklinear bersifat tak terbatas pada setiap \omterm{prop:b1:reals:intervals}{selang}
        yang taktermerosot, takkontinu di setiap titik, tak monoton pada
        \omterm{prop:b1:reals:intervals}{selang} mana pun, dan petanya atas sebarang \omterm{prop:b1:reals:intervals}{selang} bersifat \omterm{def:b1:topology:dense}{padat} di
        $\R$. Rukunkanlah dengan pertanyaan 8.
  \item Monster itu ada --- pada \omterm{def:b1:structures:subgroup}{subgrup} yang \omterm{def:b1:topology:dense}{padat}, secara konstruktif:
        pada $G = \Z + \sqrt2\,\Z$ definisikan $f(m + n\sqrt2) = n$.
        Tunjukkan bahwa $f$ terdefinisi dengan baik dan aditif pada $G$, dan bahwa
        $f$ tak terbatas pada $G \cap \intoo{0}{\varepsilon}$ untuk
        setiap $\varepsilon > 0$ \emph{(karena untuk $N$ yang tetap, hanya
        berhingga banyak $g = m + n\sqrt2 \in \intoo{0}{1}$ yang mempunyai
        $\abs n \leq N$; padahal $G \cap \intoo{0}{\varepsilon}$
        tak hingga)}. Jelaskanlah dalam satu paragraf mengapa memperluas
        $f$ yang demikian ke seluruh $\R$ menuntut sebuah basis $\R$ sebagai
        ruang vektor atas $\Q$ (yaitu \emph{basis Hamel}), yang
        keberadaannya merupakan urusan aksioma pemilihan di luar
        jilid ini.
\end{enumerate}

\medskip
\noindent\textbf{Bagian IV --- Persamaan saudarinya.} Semua
fungsi di sini bersifat \omterm{def:b1:continuity:continuous}{kontinu}.
\begin{enumerate}[resume]
  \item Misalkan $f \colon \R \to \R$ \omterm{def:b1:continuity:continuous}{kontinu}, tak identik
        $0$, dengan $f(x + y) = f(x)f(y)$. Tunjukkan $f(x) =
        f\bigl(\frac x2\bigr)^2 \geq 0$, lalu $f > 0$ di mana-mana,
        lalu $f(x) = \eu^{cx}$ untuk suatu $c$: jadi eksponensialnya
        tepat merupakan morfisme \omterm{def:b1:continuity:continuous}{kontinu} dari $(\R, +)$ ke
        $(\R^*, \times)$.
  \item Misalkan $f \colon \intoo{0}{+\infty} \to \R$ \omterm{def:b1:continuity:continuous}{kontinu}
        dengan $f(xy) = f(x) + f(y)$. Tunjukkan $f(x) = c\ln x$
        \emph{(angkutlah lewat $\exp$)}.
  \item Misalkan $f \colon \intoo{0}{+\infty} \to \intoo{0}{+\infty}$
        \omterm{def:b1:continuity:continuous}{kontinu} dengan $f(xy) = f(x)f(y)$. Tunjukkan $f(x) = x^c$.
  \item Carilah semua $f \colon \R \to \R$ yang \omterm{def:b1:continuity:continuous}{kontinu} dengan
        \[
        f(x + y) = f(x) + f(y) + f(x)f(y) .
        \]
        \emph{(Telaahlah $h = 1 + f$; lalu tanganilah kasus yang merosot
        secara tersendiri.)}
  \item (Jajargenjang) Carilah semua $f \colon \R \to \R$ yang \omterm{def:b1:continuity:continuous}{kontinu}
        dengan $f(x + y) + f(x - y) = 2f(x) + 2f(y)$: tunjukkan bahwa $f$
        genap, $f(0) = 0$, $f(nx) = n^2 f(x)$ lewat induksi, lalu
        $f(x) = f(1)\,x^2$. (Persamaan ini merupakan sidik jari
        bentuk kuadratik --- yaitu hukum jajargenjang yang mendeteksi,
        pada jilid Tahun ke-2, norma mana yang berasal dari sebuah hasil
        kali dalam.)
\end{enumerate}

\medskip
\noindent\textbf{Bagian V --- Jensen dan kecembungan titik tengah.}
\begin{enumerate}[resume]
  \item (Persamaan Jensen) Misalkan $f \colon \R \to \R$
        \omterm{def:b1:continuity:continuous}{kontinu} dengan $f\bigl(\frac{x+y}{2}\bigr) = \frac{f(x)
        + f(y)}{2}$. Tunjukkan bahwa $g = f - f(0)$ memenuhi
        $g\bigl(\frac x2\bigr) = \frac{g(x)}{2}$, lalu simpulkan bahwa
        $g$ aditif, dan tutuplah dengan $f(x) = cx + d$.
  \item Andaikan sekarang hanya \emph{ketaksamaannya}: bahwa $f$ \omterm{def:b1:continuity:continuous}{kontinu}
        dengan
        \[
        f\Bigl(\frac{x + y}{2}\Bigr) \leq \frac{f(x) +
        f(y)}{2} \qquad (x, y \in \R).
        \]
        Buktikan lewat induksi pada $n$ bahwa untuk setiap bobot diadik
        $\lambda = \frac{k}{2^n} \in \intcc{0}{1}$:
        \[
        f\bigl(\lambda x + (1 - \lambda)y\bigr) \leq \lambda
        f(x) + (1 - \lambda) f(y) .
        \]
  \item Perluaslah lewat \omterm{def:b1:continuity:continuous}{kekontinuan} dan \omterm{def:b1:topology:dense}{kepadatan} bilangan diadik
        (\cref{exo:b1:reals:8}) ke setiap $\lambda \in
        \intcc{0}{1}$: jadi fungsi \omterm{def:b1:continuity:continuous}{kontinu} yang cembung titik tengah
        memenuhi ketaksamaan kecembungan yang penuh (yaitu gagasan
        yang ditelaah secara sistematis pada \cref{ch:b1:derivative}).
  \item Tunjukkan bahwa \omterm{def:b1:continuity:continuous}{kekontinuannya} tak dapat dilepaskan: bahwa $f$ aditif
        yang taklinear memenuhi \emph{kesamaan titik tengah} pada pertanyaan
        19 namun tak memenuhi ketaksamaan kecembungan pada \omterm{prop:b1:reals:intervals}{selang} mana pun (pertanyaan
        12). Moralnya: bahwa kecembungan titik tengah adalah sifat
        bertahap terbilang (lewat bilangan diadik), sedangkan kecembungan bersifat kontinum;
        dan kekontinuanlah jembatannya --- persis seperti pada Bagian I--II.
\end{enumerate}

\medskip
\noindent\textbf{Bagian VI --- Variasi terakhir dan sintesisnya.}
\begin{enumerate}[resume]
  \item Carilah semua $f \colon \R \to \R$ yang \omterm{def:b1:continuity:continuous}{kontinu} dengan $f(x + y)
        = f(x) + f(y) + xy$ \emph{(kurangkanlah penyelesaian
        khususnya $\frac{x^2}{2}$)}.
  \item Buktikan: bahwa jika $f \colon \R \to \R$ \omterm{def:b1:continuity:continuous}{kontinu} dan
        aditif sekadar pada sebuah \emph{\omterm{def:b1:structures:subgroup}{subgrup} yang \omterm{def:b1:topology:dense}{padat}} $G$ (yakni
        $f(g + g') = f(g) + f(g')$ untuk $g, g' \in G$), maka $f$
        aditif pada $\R$. Secara lebih umum, dua fungsi \omterm{def:b1:continuity:continuous}{kontinu}
        yang bersesuaian pada sebuah \omterm{def:b1:logic:sets}{himpunan} bagian $\R$ yang \omterm{def:b1:topology:dense}{padat} bersifat
        sama.
  \item Sintesis, satu kalimat untuk masing-masing: (i) nyatakanlah tangga
        keteraturan pertanyaan 8 dari ingatan; (ii) jelaskanlah mengapa
        ``grafiknya \omterm{def:b1:topology:dense}{padat} di bidangnya'' merupakan citra batin yang tepat
        bagi kegagalan keteraturannya; (iii) daftarkanlah kelima
        fungsi klasik yang dicirikan pada Bagian IV--V beserta
        satu metode tunggal yang menangkap semuanya; (iv) sebutkanlah kedua
        tempat yang di situ \omterm{def:b1:topology:dense}{kepadatan} $\Q$ (atau bilangan diadik) di $\R$
        memikul argumennya, dan tempat yang di situ ia tak dapat
        (pertanyaan 13).

\end{enumerate}
\end{problem}
